%====================================================
% 17 Discussion
%====================================================

\section{Discussion}

In this study, five different Convolutional Neural Network (CNN) architectures have been designed and tested with MNIST dataset for handwritten digit recognition with promising results. The model modifications were made on a single CNN parameter and the rest of the network remained the same for each model. This method allowed the systematic study of the impact of architectural changes on classification accuracy.

The experimental outcomes showed that all the five CNN models successfully realized the handwritten digit recognition, and the recognition accuracy of the models exceeded 97\% in the test. The model with the highest testing accuracy was Model 3 with an additional convolution layer that resulted in a testing accuracy of 99.11% and the least number of misclassified images. The added convolution layer was able to capture higher order image features, which led to better classification performance.

When comparing the kernel sizes a change from $3\times3$ to $5\times5$ convolution kernel resulted in a marginal difference in classification accuracy. The former because of the smaller number of trainable parameters and the training time, but the latter, the smaller convolution kernel showed slightly better recognition accuracy, suggesting that the smaller convolution kernel is suitable for extracting local features from handwritten digit images.

The optimizer comparison demonstrated that the Adam optimizer was much better than the SGD optimizer. The accuracy on the test set was slightly lower for SGD, testing loss was higher, and the number of misclassifications was significantly higher. Adam's adaptive learning strategy helped to achieve more efficient optimization and faster convergence during training.

Both ReLU and Tanh were capable of excellent classification performance as pointed out by the comparison of activation functions. The model with the Tanh activation function (Model 5) obtained the lowest testing loss but still retained the same testing accuracy as the baseline model. Yet due to its simplicity of computation and efficient propagation of gradients, ReLU is still a widely used activation function for numerous deep learning applications.

The confusion matrix and Misclassification analysis also validated the effectiveness of the developed CNN models. A high percentage of the digits in the handwritten test set were successfully classified, and only a few prediction errors were made between digits that are hard to tell apart visually. The major reason for these misclassifications was due to the ambiguity in the handwriting style and not due to the CNN model.

It was observed that the trained CNN worked well beyond the original MNIST training data set, with the successful prediction of custom handwritten digit images. Several steps in the preprocessing pipeline, such as converting to grayscale, inverting, thresholds, resizing, padding, and normalization, helped the custom images to be very similar to the MNIST image format, making the CNN able to classify them correctly.

Overall, the experimental results show the highly efficient solution of CNN for handwritten digit recognition. The comparative analysis also underscores the significance of making correct architectural choices, such as the number of convolutional layers, the optimizer, the kernel size and the activation function, in order to maximize the classification accuracy, reduce the computational load and minimize the complexity of the model.